\chapter{Scope, Motivation and Project Status}

The RADOME project proposes a distributed passive electromagnetic sensing network based on geodesic and conformal radomes installed at elevated sites. Each station combines multiband receiving apertures, local digitization, event-oriented processing, precise timing, calibration and optical networking. Several stations cooperate to estimate angle of arrival (AOA), time difference of arrival (TDOA), frequency difference of arrival (FDOA), Doppler and polarization.

The term \emph{radome} derives from \emph{radar dome}: an enclosure placed over an antenna to protect it from wind, temperature, dust, rain and snow while seeking electromagnetic transparency \cite{qamar2020}. In this project, however, the dome does not house a dedicated radar transmitter. It protects and integrates receiving apertures only. The radar function emerges from cooperation among stations: each node observes direct emitter signals or signals from external illuminators reflected by targets, and the network estimates position and motion through AOA triangulation and TDOA/FDOA/Doppler multilateration and fusion. The concept is therefore a passive multistatic radar built from receiving radomes, not a set of distributed active radars.

Receive-only operation makes the station substantially less observable in the emission domain than an active radar: it generates no illumination pulses, RF power-amplifier output or scanning transmission that discloses its position and operating mode. The system listens to external waveforms of known origin and geometry; any RF energy attributable to its presence by an external sensor is secondary scattering---an echo of ambient illumination---rather than an interrogation emission generated by the radome. This electromagnetic discretion is an architectural property. It is distinct from the structure's radar cross section, whose magnitude depends on materials and geometry.

Many conventional radomes are developed for active radars, for which the enclosure must preserve both the transmitted wave and the received echo. Wide coverage, scan angle and polarization require co-design of permittivity, layer thicknesses, cores and multilayer stacks, together with specialized electromagnetic manufacturing and validation \cite{qamar2020,ramanamurthy2016}. RADOME investigates a different cost--performance balance: the receiving Yagis are exposed and oriented by each face normal, so their useful path does not depend on multiband shell transparency. This may simplify module-level characterization and replacement, but transfers requirements to exposed supports, sealing, cables, corrosion, wind, lightning and recalibration. Lower acquisition or life-cycle cost is therefore a hypothesis to be quantified, not a demonstrated result.

The network may exploit known opportunistic illuminators, including cellular infrastructure, broadcast transmitters, satellites and other authorized RF sources. In bistatic operation, the direct reference path is compared with the path involving a target and a receiver. The system is conceptual: performance, coverage, locations and accuracy must be demonstrated experimentally and are not operational guarantees.

The project objectives are to develop a measurable multiband architecture, validate a three-node VHF/UHF demonstrator, characterize end-to-end timing and RF delays, preserve event windows locally, and integrate the electromagnetic platform with resilient energy, thermal, communications and site infrastructure.

The independent English edition, assembled from \texttt{radome-en.tex}, is one of the two authoritative technical publications. ``Document version 1.1'' identifies the controlled publication, while ``architecture revision 3'' identifies the current conceptual system architecture; these identifiers are independent. Earlier Markdown, PDF and standalone LaTeX documents are historical inputs unless a current chapter explicitly incorporates their content. Quantitative values are controlled in \texttt{PARAMETERS.md}, with status proposed, simulated, measured or derived, and architecture decisions are recorded in \texttt{DECISIONS.md}. In case of conflict, an approved parameter-register entry and this article take precedence over historical material.
